% Identity Bedrock — Source LaTeX (version française)
% SAUVEGARDER SOUS: identity_bedrock_fr.tex
% Compilation : pdflatex -> bibtex -> pdflatex -> pdflatex

\documentclass[11pt,a4paper]{article}

\usepackage[utf8]{inputenc}
\usepackage[T1]{fontenc}
\usepackage[french]{babel}
\usepackage{lmodern}
\usepackage{microtype}
\usepackage{amsmath,amssymb,amsthm}
\usepackage{mathtools}
\usepackage{booktabs}
\usepackage{array}
\usepackage[colorlinks=true,linkcolor=black,citecolor=black,urlcolor=blue]{hyperref}
\usepackage{xcolor}
\usepackage{listings}
\usepackage{geometry}
\usepackage{fancyhdr}
\usepackage{url}
\usepackage{cite}
\usepackage{graphicx}

\geometry{a4paper,top=25mm,bottom=25mm,left=25mm,right=25mm}

\newtheorem{theorem}{Théorème}[section]
\newtheorem{lemma}[theorem]{Lemme}
\newtheorem{proposition}[theorem]{Proposition}
\newtheorem{corollary}[theorem]{Corollaire}
\theoremstyle{definition}
\newtheorem{definition}[theorem]{Définition}
\newtheorem{remark}[theorem]{Remarque}

\lstset{basicstyle=\small\ttfamily,breaklines=true,frame=single,
        numbers=none,xleftmargin=1em,language={}}

\pagestyle{fancy}
\fancyhf{}
\lhead{\small Identity Bedrock}
\rhead{\small Bouchard, 2026}
\cfoot{\thepage}
\renewcommand{\headrulewidth}{0.4pt}

\newcommand{\OWF}{\mathrm{FSU}}
\newcommand{\SIG}{\mathrm{SIG}}
\newcommand{\EUFCMA}{\mathrm{EUF\text{-}CMA}}
\newcommand{\negl}{\mathrm{negl}}
\newcommand{\PBAT}{\textsf{PBAT}}
\newcommand{\BD}{\textsf{Identity Bedrock}}

\begin{document}

\title{\textbf{Identity Bedrock~: opérationnalisation du substrat
cryptographique minimal de Rompel\\pour une identité souveraine avec
preuve de présence physique}}

\author{
  Jean-Fran\c{c}ois Bouchard\\[4pt]
  \small PrivateNess Network (co-fondateur, directeur technique)\\
  \small Emercoin (contributeur communautaire)\\[4pt]
  \small \texttt{[contact via privateness.network]}
}

\date{Mars 2026 \quad|\quad Version 1.0}

\maketitle

\begin{center}
\small\itshape
\textbf{Déclaration IA~:} L'auteur a utilisé un assistant de rédaction
fondé sur un modèle de langage pour la rédaction et la révision du
manuscrit. L'ensemble des contributions techniques, de la conception
du système, de l'implémentation, du cadre théorique et des intuitions
de liaison décrits dans ce travail sont ceux de l'auteur, issus de
neuf ans d'expérience opérationnelle dans le déploiement
d'infrastructures fondées sur Emercoin.
\end{center}

\vspace{4pt}\noindent\rule{\textwidth}{0.4pt}\vspace{2pt}

\begin{abstract}
Nous présentons \BD{}, un système d'identité déployé fondé
explicitement sur la preuve de Rompel selon laquelle les fonctions à
sens unique ($\OWF$) sont nécessaires et suffisantes pour les schémas
de signature existentiellement infalsifiables sous attaque à messages
choisis adaptative ($\SIG_{\EUFCMA}$)~\cite{Rompel90}. Le système
sépare la cryptographie d'identité (Ed25519 ou ML-DSA) du consensus
d'annuaire via un registre agnostique à la chaîne en ajout seul,
éliminant les clés institutionnelles sous l'individu. Déployé sur
Emercoin NVS avec ancrage Bitcoin AuxPoW, le coût d'attaque sur la
couche d'annuaire approche cinq fois la capitalisation boursière
d'EMC~\cite{Kho21} en raison de la triple barrière de consensus.

Nous identifions en outre l'audio par conduction osseuse comme un
phénomène physique satisfaisant la propriété d'entropie inaccessible
formalisée par Haitner et~al.~\cite{Haitner10}~: le désaccord
d'impédance mécanique à l'interface air-os ($T \approx
1{,}1 \times 10^{-3}$) rend le canal d'authentification
inconditionnellement inaccessible à la rejouabilité, indépendamment
de toute hypothèse calculatoire. Nous intégrons ce facteur comme
quatrième facteur d'authentification et proposons le jeton
d'authentification à liaison physique (\PBAT{}), une structure signée
portant à la fois une preuve biométrique de session et une preuve
cryptographique d'identité.

La contribution réside dans la liaison de primitives existantes en
un système dont les propriétés de sécurité émergentes n'existent dans
aucune d'elles séparément.
\end{abstract}

\medskip
\noindent\textbf{Mots-clés~:} fonctions à sens unique, signatures
numériques, identité souveraine, authentification par conduction
osseuse, entropie inaccessible, identité sur chaîne de blocs,
cryptographie post-quantique, preuve de présence physique.

\noindent\rule{\textwidth}{0.4pt}
\tableofcontents
\newpage

\section{Introduction}

\subsection{Le problème~: la capture institutionnelle}

Tout système d'identité numérique déployé place des clés
institutionnelles sous l'utilisateur~: autorités de certification,
fournisseurs d'identité, registres de clés, opérateurs de plateforme.
Ces entités peuvent révoquer, refuser, surveiller ou contrôler
l'identité sans le consentement de l'utilisateur.

Le modèle de menace motivant ce travail suppose un adversaire
disposant de ressources conventionnelles substantielles, d'un soutien
politique potentiel et d'un contrôle total de l'infrastructure
d'identité traditionnelle. La question centrale est~:
\emph{peut-il exister une identité sans clés institutionnelles sous
l'individu~?}

\subsection{La réponse théorique~: Rompel 1990}

Rompel a démontré que les $\OWF$s sont à la fois nécessaires et
suffisantes pour les schémas de signature
$\SIG_{\EUFCMA}$~\cite{Rompel90}. Il s'agit du minimum démontrable
--- rien de strictement plus faible ne soutient l'authentification
publique sous hypothèses standard.

Notre thèse~: fixer la $\OWF$ comme plancher cryptographique et y
ajouter uniquement un annuaire agnostique à la chaîne en ajout seul
produit une construction d'identité qui minimise à zéro les points
de capture institutionnelle sous la clé individuelle.

\subsection{L'intuition de liaison}

Ce travail suit le schéma des innovations par liaison en cryptographie
appliquée~: la valeur ne réside pas dans l'invention de nouvelles
primitives, mais dans la reconnaissance que des primitives existantes,
liées d'une façon spécifique, produisent des propriétés émergentes
qu'aucune d'elles ne présente seule.

\subsection{Contributions}

\paragraph{Contribution 1 --- Théorie.}
Formalisation d'un système d'identité agnostique à la chaîne fondé
sur Rompel.

\paragraph{Contribution 2 --- Systèmes.}
Architecture~: couche d'identité Ed25519 ou ML-DSA~\cite{FIPS204}
avec substrat d'annuaire interchangeable. Implémentation~: Emercoin
NVS~\cite{Emercoin} et NESS Fiber~\cite{DPO}. Déployé et vérifiable
publiquement depuis 2013.

\paragraph{Contribution 3 --- Extension de présence physique.}
Identification de l'audio par conduction osseuse comme phénomène
physique satisfaisant la propriété d'entropie inaccessible de Haitner
et~al.~\cite{Haitner10}~\cite{HeadSonic25,HCRAuth24}. Spécification
du jeton \PBAT{}.

\subsection{Organisation de l'article}

Section~\ref{sec:background}~: contexte. Section~\ref{sec:threat}~:
modèle de menace. Section~\ref{sec:architecture}~: architecture.
Section~\ref{sec:implementation}~: implémentation.
Section~\ref{sec:physical}~: présence physique.
Section~\ref{sec:evaluation}~: évaluation.
Section~\ref{sec:discussion}~: limitations et travaux futurs.

\section{Contexte et travaux connexes}
\label{sec:background}

\subsection{Le théorème de Rompel (1990)}

Les $\OWF$s, UOWHFs et schémas $\SIG_{\EUFCMA}$ sont équivalents en
force~\cite{NaorYung89}. Rompel~\cite{Rompel90} et Haitner
et~al.~\cite{Haitner10} ont montré comment construire des UOWHFs
depuis des $\OWF$s arbitraires via l'entropie inaccessible~\cite{KatzKoo05}.
\[
  \OWF \;\Leftrightarrow\; \SIG_{\EUFCMA}
\]
sous les définitions standard~\cite{Rompel90,Goldreich01}.

\subsection{Entropie inaccessible}
\label{sec:inaccessible}

\begin{definition}[Entropie inaccessible, informelle~{\cite{Haitner10}}]
\label{def:inaccessible}
Soit $F: \{0,1\}^n \to \{0,1\}^m$. L'\emph{entropie inaccessible}
de $F^{-1}$ est l'écart
\[
  \Delta \;=\; H(X \mid F(X)) \;-\; \max_{A \in \mathrm{PPT}}
  H(A(X;\,R) \mid X)
\]
où $X$ est uniforme sur $\{0,1\}^n$ et $A$ est contraint à produire
des éléments de $F^{-1}(F(X))$.
\end{definition}

\subsection{Systèmes d'identité fondés sur la chaîne de blocs}

\textbf{Namecoin / EmerDNS.} Lacune~: clés d'identité liées aux
clés de portefeuille~; mineurs comme AC implicite.

\textbf{DID W3C / Ethereum.} Lacune~: clés DAO, administrateurs de
registre, dépendances fournisseurs.

\textbf{Sovrin / Hyperledger Indy.} Lacune~: clés d'intendants
contrôlent le réseau.

\textbf{Distinction de \BD{}.} Plancher $\OWF$ pur. Séparation
agnostique à la chaîne. Aucune clé institutionnelle sous
l'utilisateur.

\subsection{AuxPoW Bitcoin et minage fusionné}

Khovayko~\cite{Kho21} introduit le coefficient \emph{hack51}
($\text{coût\_attaque} / \text{capitalisation}$). Le triple consensus
d'Emercoin (PoS + PoW + minage fusionné) requiert des ressources
approchant cinq fois sa capitalisation boursière.

\subsection{Authentification biométrique par conduction osseuse}

HeadSonic~\cite{HeadSonic25}~: 96,59\% de précision, résistance au
rejeu empiriquement démontrée. HCR-Auth~\cite{HCRAuth24}~: idem.
Aucun de ces travaux ne fournit l'argument formel. Nous le fournissons
à la section~\ref{sec:impedance}.

\subsection{Consensus Skycoin Fiber}

Skycoin Fiber~\cite{Skycoin}~: clés d'éditeur annoncées. Tout nombre
de n\oe{}uds hostiles jusqu'à $N-1$ ne peut pas forger de blocs sans
ces clés. Le moteur BoltDB en ajout seul rend la réécriture de l'historique impossible~: la seule attaque disponible à un adversaire détenant une clé d'éditeur volée est une réinitialisation au bloc de genèse~0, elle-même maximalement visible par tous les participants.

\subsection{Synthèse des différences}

\BD{} se distingue sur trois points~: plancher Rompel explicite~;
architecture agnostique à la chaîne~; premier argument formel pour
la résistance au rejeu des biométriques par conduction osseuse. Le
\PBAT{} est le premier format de jeton portant preuve biométrique et
cryptographique dans une structure signée unique.

\section{Modèle de menace et exigences}
\label{sec:threat}

\subsection{Capacités de l'adversaire}

\textbf{Supposées~:} ressources substantielles~; contrôle de
l'infrastructure d'identité~; adversaire classique (temps polynomial).

\textbf{Non supposées~:} ne peut pas casser la $\OWF$~; ne peut pas
mener une attaque à 51\% soutenue sur Bitcoin~; ne peut pas violer
la physique acoustique.

\textbf{Périmètre~:} adversaires classiques. Migration ML-DSA
discutée à la section~\ref{sec:limitations}.

\subsection{Exigences de conception}

\textbf{R1~:} $\OWF$ comme plancher. \textbf{R2~:} Aucune clé
institutionnelle. \textbf{R3~:} Agnosticisme à la chaîne.
\textbf{R4~:} Coûts quantifiables. \textbf{R5~:} Vérification
opérationnelle. \textbf{R6~:} Présence physique inconditionnelle.

\section{Architecture d'Identity Bedrock}
\label{sec:architecture}

\subsection{La conception à deux couches}

\paragraph{Couche 1~: cryptographie d'identité.}
\[
\text{Graine}\;(\geq 256\text{ bits})
\;\xrightarrow{\mathsf{GénClé}}\;(\mathit{sk},\;\mathit{pk})
\]
Ed25519~\cite{RFC8032} par défaut. ML-DSA-65~\cite{FIPS204} comme
cible de migration.

\paragraph{Couche 2~: substrat d'annuaire (interchangeable).}

\begin{definition}[Annuaire en ajout seul avec finalité]
\label{def:directory}
Un annuaire $\mathcal{D} \subseteq \mathcal{ID} \times \mathcal{PK}$
est \emph{en ajout seul avec finalité} si~: \emph{(i)} une fois que
$\mathsf{Lier}(\mathit{id}, \mathit{pk}) = 1$ est confirmé, cela ne
peut être mis à $0$ sans attaque de coût quantifiable~; et
\emph{(ii)} la liaison est globalement vérifiable sans intermédiaire.
\end{definition}

\subsection{Protocole d'authentification}

\paragraph{Configuration.}
\begin{enumerate}
  \item Génération de $(\mathit{sk}, \mathit{pk})$.
  \item Publication de $(\mathit{id}, \mathit{pk})$ dans $\mathcal{D}$.
  \item Attente de finalité.
  \item $\mathsf{Lier}(\mathit{id}, \mathit{pk}) = 1$.
\end{enumerate}

\paragraph{Défi-réponse.}
\begin{enumerate}
  \item $V \to P$~: nonce frais $c \xleftarrow{\$} \{0,1\}^\lambda$.
  \item $P \to V$~: $\sigma = \mathsf{Sign}_{\mathit{sk}}(c)$.
  \item $V$ vérifie liaison et signature~; accepte si les deux tiennent.
\end{enumerate}

\subsection{Preuve de sécurité}

\begin{theorem}
\label{thm:unforgeability}
Si le schéma est $\SIG_{\EUFCMA}$ sûr et $\mathcal{D}$ en ajout seul
avec finalité, alors \BD{} fournit une authentification infalsifiable
contre tout adversaire classique.
\end{theorem}

\begin{proof}
Le faussaire $\mathcal{B}$ simule $\mathcal{D}$ localement, transmet
le défi $c^*$ du jeu $\SIG_{\EUFCMA}$ à $\mathcal{A}$, et utilise
la réponse $\sigma^*$ comme falsification existentielle. Par sécurité
du schéma~: $\Pr[\mathcal{A}\text{ gagne}] \leq \negl(\lambda)$.\hfill\qed
\end{proof}

\section{Implémentation~: Emercoin + NESS Fiber}
\label{sec:implementation}

L'architecture décrite dans ce travail s'appuie sur la conception
publiée dans le LitePaper PrivateNess Network d'octobre
2020~\cite{Litepaper20}, qui établissait l'architecture double-aveugle
Emercoin + Skywire, la restriction des communications à 127.0.0.1,
l'entropie UHE-PRNG à 1536 bits d'état interne, et l'approche de
décentralisation de la pile OSI complète. Identity Bedrock formalise
les fondements cryptographiques de ce système déployé.

\subsection{Annuaire~: Emercoin NVS}

Consensus~: 85\% PoS + 15\% Bitcoin AuxPoW~\cite{Emercoin}.
Finalité~: $\approx$2 minutes. Liaisons ancrées via le schéma WORM
(World Object MaRkup language)~\cite{WORM}, conçu par Aleksej
Sokolov~:

\begin{lstlisting}[language=XML]
<worm>
  <user type="ness" key="<pubkey_hex>" nonce="<nonce>">
    <keys>
      <key hash="<hash>" algo="ed25519"><pubkey_hex></key>
    </keys>
    <node type="ness" url="<skywire_pubkey>"
          nonce="<nonce>" tags="enregistreur,voix"/>
  </user>
</worm>
\end{lstlisting}

\subsection{Infrastructure~: NESS Fiber}

Clés d'éditeur publiées dans \texttt{dpo/PrivateNESS.Network}~\cite{DPO}.
La compromission d'une clé est maximalement détectable~: une clé d'éditeur volée ne peut pas réécrire les blocs passés (BoltDB en ajout seul rend la réécriture impossible). La seule action disponible est une réinitialisation au bloc de genèse~0, immédiatement visible par tous les participants. La rotation des clés est elle-même un événement publiquement auditable.

\subsection{Preuve d'auto-dépense canonique}

Liée à l'identité Ed25519, ancrée en NVS. Instance déployée~:
bloc Emercoin~\#445669, 1\textsuperscript{er} novembre 2020,
22~h~16~min~45~s~UTC.
\[
\text{Bitcoin AuxPoW}
\to \text{Bloc \#445669}
\to \text{DPO}
\to \text{BlockchainPubkeyStr}
\to \text{GenesisSignatureStr}
\to \text{Chaîne NESS Fiber}
\]

\subsection{Génération de clés et entropie}

Génération locale via UHEPRNG~\cite{Gibson} avec entrée d'entropie
\texttt{os.urandom} (plusieurs blocs de 1024 octets), et sources
locales optionnelles (gigue CPU, identifiants matériels). Plancher~:
$\approx$128 bits de sécurité classique.

\section{Extension de présence physique}
\label{sec:physical}

\subsection{Motivation}

Couche supplémentaire garantissant inconditionnellement la présence
physique, non par détection algorithmique de vivacité.

\subsection{Le problème de symétrie de canal}

Canal d'authentification et canal de rejeu acoustiques sont
physiquement identiques~: vulnerabilité fondamentale irréductible
par logiciel.

\subsection{La conduction osseuse brise la symétrie}
\label{sec:impedance}

$Z_\text{air} \approx 415$~Pa$\cdot$s/m,
$Z_\text{os} \approx 1{,}5 \times 10^6$~Pa$\cdot$s/m~\cite{Acoustics}.

\begin{equation}
\label{eq:transmission}
T = \frac{4\, Z_\text{air}\, Z_\text{os}}{(Z_\text{air}+Z_\text{os})^2}
\approx 1{,}1 \times 10^{-3}
\end{equation}

\begin{remark}
L'équation~(\ref{eq:transmission}) est une conséquence de la physique
acoustique, non une hypothèse de dureté. Elle tient inconditionnellement.
\end{remark}

\subsection{Entropie inaccessible physique}

\begin{remark}[Instanciation physique de l'entropie inaccessible]
\label{rem:physical_ie}
Le canal de conduction osseuse produit une distribution dont l'entropie
est inaccessible à tout adversaire sans contrôle physique du corps
inscrit. L'inaccessibilité est imposée par
l'équation~(\ref{eq:transmission}), non par hypothèse calculatoire~:
garantie strictement plus forte que la Définition~\ref{def:inaccessible}.
\end{remark}

\begin{remark}[Périmètre]
Observation et analogie, non théorème achevé. Réduction formelle~:
travail futur.
\end{remark}

\subsection{Pile à quatre facteurs}

\begin{table}[h]
\centering
\caption{Pile d'authentification à quatre facteurs}
\label{tab:fourfactor}
\begin{tabular}{llll}
\toprule
Facteur & Mécanisme & Attaque & Garantie \\
\midrule
Ed25519 / $\OWF$ & Plancher Rompel & Falsification & Calculatoire \\
EmerNVS / AuxPoW & Annuaire ancré & Réécriture & Économique \\
Conduction osseuse & Couplage physique & Rejeu & Inconditionnelle \\
Empreinte + \PBAT{} & Session signée & Usurpation & Biom.+crypto \\
\bottomrule
\end{tabular}
\end{table}

\subsection{Jeton d'authentification à liaison physique}

\begin{definition}[\PBAT{}]
\label{def:pbat}
Structure signée contenant~: clé publique d'identité~; preuve de
liaison d'annuaire~; empreinte biométrique de session capturée via
canal physiquement inaccessible~; nonce à usage unique~; horodatage~;
signature par la clé privée d'identité.
\end{definition}

\begin{lstlisting}
{
  "identity_pubkey":   "<pubkey Ed25519 ou ML-DSA>",
  "identity_algo":     "ed25519 | mldsa65",
  "nvs_record":        "<txid WORM>",
  "voiceprint_hash":   "<SHA256 ou SHA3-256>",
  "session_embedding": "<empreinte session courante>",
  "similarity_score":  0.94,
  "capture_channel":   "bone_conduction_msbc",
  "headset_mac":       "<adresse MAC BT>",
  "timestamp":         "<ISO8601 UTC>",
  "nonce":             "<256 bits>",
  "expires_seconds":   300,
  "signature":         "<Sign_sk(tous les champs)>"
}
\end{lstlisting}

\paragraph{Doctrine de présence physique.}
\emph{Le dispositif authentifie le corps, non les justificatifs.
Aucun chemin réseau ne substitue la présence physique.}

\section{Évaluation}
\label{sec:evaluation}

\subsection{Analyse des coûts d'attaque}
\label{sec:eval_attack}

\begin{table}[h]
\centering
\caption{Analyse des coûts d'attaque}
\label{tab:attacks}
\renewcommand{\arraystretch}{1.2}
\begin{tabular}{p{2.8cm}p{3.5cm}p{3cm}p{3cm}}
\toprule
Vecteur & Prérequis & Coût & Résultat \\
\midrule
Casser la signature & Violer $\OWF$ & Infaisable & Auth.\ s'effondre \\
Réécrire NVS & PoS+PoW+AuxPoW 51\% & $\approx\!5\times$ cap.\ EMC~\cite{Kho21} & Réécriture \\
Compromettre Fiber & Voler clés éditeur & Élevé & DoS uniquement \\
Rejouer conduction & Reproduire sig.\ mécanique & Infaisable ($T\approx 10^{-3}$) & Couche biom.\ tient \\
Contraindre user & Physique & Variable & Hors périmètre \\
\bottomrule
\end{tabular}
\end{table}

\subsection{Comparaison}

\begin{table}[h]
\centering
\caption{Comparaison des systèmes d'identité}
\label{tab:comparison}
\begin{tabular}{lcccc}
\toprule
Système & Plancher & Sans inst. & Agn.\ chaîne & Présence phys. \\
\midrule
\textbf{\BD{}} & Ed25519/Rompel & \checkmark & \checkmark & Physique (os) \\
PKI X.509 & RSA/ECDSA & $\times$ & $\times$ & Aucune \\
Namecoin & ECDSA & $\times$ & $\times$ & Aucune \\
DID Ethereum & Variable & $\times$ & $\times$ & Aucune \\
Sovrin & Ed25519 & $\times$ & $\times$ & Aucune \\
FIDO2 & ECDSA & Partiel & $\times$ & Jeton mat. \\
HeadSonic~\cite{HeadSonic25} & Aucun & N/A & N/A & Physique (os) \\
\bottomrule
\end{tabular}
\end{table}

\subsection{Performances}

\begin{itemize}
  \item Finalité NVS~: $\approx$2 minutes
  \item Défi-réponse~: millisecondes
  \item Vérification empreinte (resemblyzer\footnote{%
        \url{https://github.com/resemble-ai/resemblyzer}},
        RK3528)~: ${<}1$ seconde
  \item Génération \PBAT{}~: ${<}10$ ms
\end{itemize}

\section{Discussion}
\label{sec:discussion}

\subsection{Limitations}
\label{sec:limitations}

\paragraph{Hors périmètre.} Compromission par logiciel malveillant,
ingénierie sociale~: orthogonaux au modèle.

\paragraph{Adversaire quantique.} Migration vers ML-DSA-65
(FIPS~204~\cite{FIPS204})~: remplacement compatible à l'interface
gén-clés/signer/vérifier. Schéma WORM \texttt{<keys>} supporte
nativement plusieurs algorithmes. Couche physique inchangée.
OpenSSL~3.5 (avril~2025)~: ML-DSA natif dans Armbian.

\paragraph{Réduction formelle.} La Remarque~\ref{rem:physical_ie}
est une observation. Réduction formelle~: travail futur.

\paragraph{Contrainte.} Surface résiduelle~: contrainte,
compromission d'inscription, cryptanalyse post-quantique.

\subsection{Sur la nature de la contribution}

Liaison de primitives existantes produisant des propriétés émergentes
qu'aucune n'exhibe seule. Système opérationnel depuis
2013~\cite{Emercoin} / 2020~\cite{DPO}.

\subsection{Travaux futurs}

\begin{itemize}
  \item Réduction formelle pour la Remarque~\ref{rem:physical_ie}
  \item Preuve de sécurité complète du \PBAT{}
  \item Signatures à seuil ($k$-sur-$n$ ML-DSA)
  \item Portabilité agnostique~: Bitcoin OP\_RETURN, Ethereum
  \item Études utilisateurs~: inscription et taux de rejet erroné
\end{itemize}

\section*{Conclusion}
\addcontentsline{toc}{section}{Conclusion}

\BD{} atteint son objectif principal --- absence de clés
institutionnelles sous l'individu --- via une architecture agnostique
à la chaîne. L'entropie inaccessible physique de la conduction
osseuse (Remarque~\ref{rem:physical_ie}) produit un système à quatre
facteurs (Tableau~\ref{tab:fourfactor}) où chaque attaque est traitée
par un mécanisme catégoriquement distinct.

L'enregistrement canonique~\cite{DPO} constitue une preuve
opérationnelle~: ancré dans Bitcoin, auditable par toute partie
sans intermédiaire. Ce n'est pas un article décrivant ce qui
pourrait être construit. C'est un article décrivant ce qui fonctionne
depuis 2020.

\section*{Remerciements}
\addcontentsline{toc}{section}{Remerciements}

L'auteur remercie Oleg Khovayko pour l'infrastructure Emercoin et
le cadre \emph{hack51}~; Aleksej Sokolov (Alex014) pour la
conception du format WORM et l'implémentation de
\texttt{privatenesstools} et \texttt{pyuheprng} à partir des
spécifications de l'auteur~; Brandon <<~Synth~>> Smietana et Moses
Narrow (\texttt{moe\_narrow} / \texttt{0pcom}), responsable du
packaging Linux de Skywire~; ainsi que la communauté PrivateNess
Network depuis le 1\textsuperscript{er} septembre 2020~\cite{BTCTalk20}.
L'engagement de l'auteur dans l'infrastructure Emercoin s'étend sur
neuf ans.

\bibliographystyle{plain}
\begin{thebibliography}{99}

\bibitem{Rompel90}
J.~Rompel.
\newblock One-way functions are necessary and sufficient for secure
  signatures.
\newblock In {\em Proc.\ 22nd Annual ACM STOC}, pages 387--394, 1990.

\bibitem{Haitner10}
I.~Haitner, T.~Holenstein, O.~Reingold, S.~Vadhan, and H.~Wee.
\newblock Universal one-way hash functions via inaccessible entropy.
\newblock In {\em EUROCRYPT~2010}, LNCS vol.~6110, pages 616--637.
  Springer, 2010.

\bibitem{RFC8032}
S.~Josefsson and I.~Liusvaara.
\newblock {EdDSA}. RFC 8032, IETF, 2017.
\newblock \url{https://www.rfc-editor.org/rfc/rfc8032}

\bibitem{Gibson}
S.~Gibson.
\newblock {UHEPRNG}. Gibson Research Corporation.
\newblock \url{https://www.grc.com/otg/uheprng.htm}

\bibitem{Emercoin}
Emercoin Development Team.
\newblock Emercoin. \url{https://emercoin.com}, 2013--2026.

\bibitem{Kho21}
O.~Khovayko.
\newblock {Оценка стоимости атаки 51\% на Emercoin: hack51}.
\newblock Dzen, 2021.
\newblock \url{https://dzen.ru/a/YUotZkN9nAeElPEQ} [En russe]

\bibitem{CoinMetrics24}
L.~Nuzzi et~al.
\newblock Breaking {BFT}. Coin Metrics, 2024.

\bibitem{Skycoin}
B.~Smietana (<<~Synth~>>) et~al.
\newblock Skycoin. \url{https://www.skycoin.com}, 2014--2026.
\newblock Packaging~: M.~Narrow (\texttt{0pcom}).

\bibitem{WORM}
A.~Sokolov (Alex014) / {NESS-Network}.
\newblock {WORM} --- {W}orld {O}bject {Ma}R\-kup language.
\newblock \url{https://github.com/ness-network/worm}

\bibitem{NaorYung89}
M.~Naor and M.~Yung.
\newblock Universal one-way hash functions.
\newblock In {\em Proc.\ 21st Annual ACM STOC}, pages 33--43, 1989.

\bibitem{KatzKoo05}
J.~Katz and C.-Y. Koo.
\newblock ePrint 2005/328, 2005.
\newblock \url{https://eprint.iacr.org/2005/328}

\bibitem{FIPS204}
{NIST}.
\newblock {ML-DSA}. {FIPS}~204, août 2024.
\newblock \url{https://doi.org/10.6028/NIST.FIPS.204}

\bibitem{HeadSonic25}
Z.~He et~al.
\newblock {HeadSonic}.
\newblock {\em IEEE Trans.\ Mobile Comput.}, 24(9):7914--7928, 2025.
\newblock doi:10.1109/TMC.2025.3551272

\bibitem{HCRAuth24}
Z.~He et~al.
\newblock {HCR-Auth}.
\newblock {\em Proc.\ ACM IMWUT}, 8(4):1--27, 2024.
\newblock doi:10.1145/3699780

\bibitem{Acoustics}
L.~E.~Kinsler et~al.
\newblock {\em Fundamentals of Acoustics}, 4th~ed.
\newblock Wiley, 2000. ISBN~978-0-471-84789-2.

\bibitem{Goldreich01}
O.~Goldreich.
\newblock {\em Foundations of Cryptography, Vol.~1}.
\newblock Cambridge University Press, 2001.

\bibitem{BTCTalk20}
{PrivateNESS Network}.
\newblock Bitcointalk, 1\textsuperscript{er} septembre 2020.
\newblock \url{https://bitcointalk.org/index.php?topic=5272741.0}

\bibitem{Litepaper20}
J.-F.~Bouchard et~al.
\newblock {PrivateNess Network LitePaper}.
\newblock PrivateNess Network, 9~octobre 2020.
\newblock \url{https://privateness.network/?p=255}
\newblock \textit{<<~La seule arme efficace contre la big tech,
  c'est la small tech décentralisée.~>>}

\bibitem{DPO}
{PrivateNESS Network}.
\newblock EmerNVS DPO, bloc~\#445669, novembre 2020.
\newblock \url{https://explorer.emercoin.com/dpo/PrivateNESS.Network}

\end{thebibliography}

\appendix

\section{Définitions formelles}
\label{app:defs}

\begin{definition}[Fonction à sens unique]
$f$ est une FSU si calculable en temps polynomial et
$\Pr[\mathcal{A}(f(x)) \in f^{-1}(f(x))] \leq \negl(\lambda)$
pour tout adversaire PPT $\mathcal{A}$.
\end{definition}

\begin{definition}[$\SIG_{\EUFCMA}$]
$(\mathsf{Gén},\mathsf{Sign},\mathsf{Vérif})$ est EUF-CMA sûr si
$\Pr[\mathsf{Vérif}(pk,m^*,\sigma^*)=1 \wedge m^*\notin Q]
\leq \negl(\lambda)$.
\end{definition}

\begin{definition}[Canal biométrique physiquement inaccessible]
\label{def:physical_channel}
Canal dont l'entropie est inaccessible sans contrôle physique du sujet
inscrit, par propriétés physiques du canal non par hypothèse
calculatoire.
\end{definition}

\section{Schéma WORM}
\label{app:worm}

\begin{lstlisting}[language=XML]
<worm>
  <user type="ness" key="<pubkey>" nonce="<nonce>">
    <keys>
      <key hash="<h>" algo="ed25519"><pubkey></key>
      <key hash="<h>" algo="mldsa65"><mldsa_pubkey></key>
    </keys>
    <node type="ness" url="<skywire_pk>"
          nonce="<n>" tags="enregistreur,voix"/>
  </user>
</worm>
\end{lstlisting}

\section{Vérification du déploiement}
\label{app:deployment}

\begin{itemize}
  \item DPO~: \url{https://explorer.emercoin.com/dpo/PrivateNESS.Network}
        (bloc~\#445669, 1\textsuperscript{er} nov.\ 2020, 22~h~16~UTC)
  \item Bitcointalk~: \url{https://bitcointalk.org/index.php?topic=5272741.0}
  \item \url{https://privateness.network}
  \item \url{https://github.com/ness-network/privatenesstools}
  \item \url{https://github.com/ness-network/pyuheprng}
  \item \url{https://github.com/ness-network/worm}
  \item \url{https://ness-explorer.magnetosphere.net}
  \item \url{https://explorer.privateness.network}
\end{itemize}

\end{document}